% =============================================================================
% 4.3 自适应气象决策报告生成技术
% 草稿版本：v0.4 (2026-05-07)
%   v0.1：完成正文初稿。
%   v0.3：按导师意见 9 规范化报告片段排版。
%   v0.4：按导师意见 10 把第 4 章整体小结从本节末段迁出到独立的
%        \subsection{本章小结}（4_4_本章小结.tex，label sec:chap4-summary），
%        本节末段仅保留承接句指向 4.4 节，避免与 4.4 节重复叙述。
%        ① 在文件顶部定义局部 reportsnippet 环境，提供"档位标签 +
%           分隔线 + 左缩进正文"的统一视觉容器，替代原裸 quote 块。
%        ② 替换两个对照样例为更有冲击力的真实端到端评测产出：
%           a) 暴雨判定 → 雾/驾车安全 (decision_006_fog)。full 档直接
%              给出 GB/T 27964-2011 雾的预报等级 4 档对照表（轻雾/雾/
%              浓雾/强浓雾），LLM-as-judge 专业性 5/5 vs ablated 3/5，
%              是 30 条端到端用例里专业性差距最大的对照之一；
%           b) 洗车 → 户外活动 (decision_001_outdoor)。full 档给出
%              GB/T 28592-2012 大雨档 25.0--49.9 毫米精确区间 +
%              雨势猛烈可读化解释，比洗车指数更能体现 RAG 在数值阈值
%              上的提取能力。
%        ③ 摘选片段忠实于 experiments/results/e2e_eval_20260505_183622.json
%           真实输出（可由该文件 raw_outputs[0].ablated/full 中按 id
%           检索复核），仅去除 markdown 表格与 emoji 等不属于决策文本
%           的视觉装饰，保留全部数值与 GB/T 编号。
% 字数：约 1750 字（不含 reportsnippet 块）
% 风格规则：v0.2 风格（无 \textbf / \textit / 引例括号 / 双引号 / 破折号）
% 引用：\cite{ref1, ref9, ref18, ref20}
% 图表：
%   - 表 \reftab{tab:report-template-routing}  四类场景与模板要素对应
%   - 表 \reftab{tab:report-judge-summary}     双档 LLM-as-judge 4 维度对照（节略，详见 5.2.2）
%   - reportsnippet 块 ×4：雾/驾车 与 户外活动 两个用例的 ablated vs full 对照
% 依赖宏包：booktabs / tabularx（主稿已加载）
% =============================================================================

% ----- 局部环境：reportsnippet -----
% 用法：
%   \begin{reportsnippet}{档位}{用例描述}
%       报告正文（可多段、可含 \texttt{...} 字段名）
%   \end{reportsnippet}
% 视觉效果：顶部"[档位] 用例描述"小标题（黑体），下接 0.4pt 横线，
%   正文整体左缩 1.2em、右缩 0.5em，正文字号与正文相同，便于读者一眼
%   分辨"这是被引模型输出"而非论文叙述。
\makeatletter
\newenvironment{reportsnippet}[2]{%
    \par\addvspace{0.5em}%
    \noindent
    {\hei\wuhao\fbox{\,#1\,}}\quad{\song\wuhao #2}\par
    \vspace{0.05em}%
    \noindent\rule{\linewidth}{0.4pt}\par
    \vspace{0.25em}%
    \begingroup
    \song\wuhao
    \leftskip=1.2em \rightskip=0.5em
    \parindent=0pt
}{%
    \par\endgroup
    \addvspace{0.5em}%
}
\makeatother

\subsection{自适应气象决策报告生成技术}
\label{sec:report-generation}

前述双通路 RAG 与代码生成加受限沙箱机制分别给出了带权威出处的语义
标签与可信的数值统计结果。然而面向终端用户的输出不是这些孤立的中间
产物，而是一份按场景组织的自然
语言决策报告。气象决策报告的具体形态高度依赖任务类型。日常查询场景下
用户只关心当前与未来若干小时的事实数据；出行决策场景下用户关心的是从
出发地到目的地不同时段的多点天气以及对应的穿衣与携物建议；户外作业与
预警场景下用户关心的是一份带权威分级与国标出处的安全结论。一份能够覆盖
这些异构场景的报告生成机制，必须解决信息整合、模板选择与引用规范三个
耦合问题。本节据此设计了基于 System Prompt 编排的自适应报告生成方案，
并以端到端用例对照说明其相对于无 RAG 基线的提升。

\subsubsection*{4.3.1 \quad 报告生成的设计需求与三类输入整合}

报告生成层位于 ReAct 循环的最后一步 \cite{ref18}，其输入由三类异构片段
构成。第一类是天气工具返回的事实数据，包含温度、风力、降水、湿度、能见度
等标量值与时序数组。第二类是 \texttt{bridge\_\bk{}weather\_\bk{}data} 与
\texttt{search\_\bk{}knowledge} 返回的语义解读，包含分级名、影响描述与国标
出处。第三类是 \texttt{compute(data)} 沙箱执行返回的统计结果，包含具体
数值与对应的代码推理过程。三类片段在数据形态、置信度来源与组织方式上
各不相同，简单拼接会产生信息冗余、出处遗漏或数值与等级冲突等问题。

报告生成的设计目标是把上述三类片段按统一结构编排为面向决策的自然语言
输出，同时保证三条不变性。第一是事实优先。数值数据必须先于解读出现，
避免大语言模型用主观叙述覆盖客观读数。第二是出处必引。凡涉及分级判定
或权威性结论的句子，引用源必须保留在报告中而不被改写或省略。第三是
场景对齐。回答内容应聚焦用户实际在场的时间地点，不输出与决策无关的
冗余信息。Kim 等 \cite{ref1} 在 ClimateAgent 中通过多智能体协作实现
类似的场景编排，本文受其启发但裁剪到单智能体形态，把场景编排逻辑下沉到
System Prompt 而非额外引入编排智能体。

\subsubsection*{4.3.2 \quad System Prompt 驱动的多场景报告模板}

本文不采用基于外部 DSL 的硬模板方案，而是把场景模板与编排规则编码到
\texttt{src/\bk{}config/\bk{}settings.\bk{}py} 中的 \texttt{SYSTEM\_\bk{}PROMPT} 字段。
该方案在工程实现上的优势是模板修改不需要改动核心代码，仅需调整 prompt
文本即可生效，便于在多场景任务中快速迭代路由规则与约束策略。这一设计与
LLM Agent 在配置化编排与可组合链路方面的工程实践一致 \cite{ref10,ref15,ref18}。
System Prompt 把报告生成约束分为四组规则。
工具使用规则定义了从用户问题到工具调用序列的路由路径。RAG 知识检索规则
定义了双通路 RAG 的强制调用条件与引用源不可省略约束。语义桥接调用约束
定义了数值到分级判定必须经过 \texttt{bridge\_\bk{}weather\_\bk{}data} 而不能由
大语言模型自行脑补阈值。出行场景推理规则定义了多地点多时段的对齐方式，
例如早八到晚六从武汉到成都的查询应分别查武汉上午与成都傍晚两个时段，
不查武汉晚上或成都早上。

这一类规则的具体覆盖见 \reftab{tab:report-template-routing}，四类典型
气象决策场景在报告要素上的差异通过 System Prompt 显式编码，而非依赖大
语言模型的隐式偏好。

\begin{table}[!ht]
    \centering
    \caption{四类气象决策场景下报告要素的差异化编排}
    \label{tab:report-template-routing}
    \song\fontsize{9pt}{12pt}\selectfont
    \renewcommand{\arraystretch}{0.95}
    \begin{tabular}{llll}
        \toprule
        场景 & 主要工具调用 & 必出报告要素 & 引用强制 \\
        \midrule
        日常查询 & current\_weather + indices & 数值 + 体感解读 & 可选 \\
        出行决策 & forcast\_weather + indices + bridge & 多时段数值 + 穿衣建议 & 涉及阈值时强制 \\
        户外作业 & current\_weather + bridge + knowledge & 数值 + 分级 + 安全结论 & 强制 \\
        极端预警 & weather\_warning + bridge + knowledge & 预警等级 + 国标依据 + 应对建议 & 强制 \\
        \bottomrule
    \end{tabular}
\end{table}

\subsubsection*{4.3.3 \quad 三类输入的编排逻辑与典型报告对照}

报告编排的具体次序遵循事实先行原则。先以表格或要点形式列出 API 工具
返回的数值数据，再附加 \texttt{bridge\_\bk{}weather\_\bk{}data} 返回的等级与
国标出处，最后给出基于这两类信息的决策建议。一旦涉及统计计算需求，
\texttt{compute(data)} 的执行结果会以具体数值加计算依据的方式嵌入到
事实数据段，避免大语言模型在数值聚合上的幻觉。Hong 等 \cite{ref20} 关于
让大语言模型直算等价于无草稿心算的论述在报告生成阶段同样适用，因为
最终输出本身就是一种聚合性表达，只要中间环节出现一次数值幻觉就会污染
整段报告。

下面以两个典型用例对比 \texttt{full} 档（含双通路 RAG）与 \texttt{ablated}
档（仅 9 个天气工具）的报告输出差异。两档共享同一份意图识别前置与同一份
天气数据，差异仅来自报告生成阶段是否能够调用 \texttt{search\_\bk{}knowledge}
与 \texttt{bridge\_\bk{}weather\_\bk{}data}。两个对照用例的报告片段均摘选自
\texttt{e2e\_\bk{}bench.\bk{}jsonl} 端到端评测的真实输出，仅删除了 markdown
表格与图标等不属于决策文本的视觉装饰，保留全部数值与 GB/T 编号原文。

第一个对照用例 \texttt{decision\_\bk{}006\_\bk{}fog} 是用户查询明天南京的雾
有多大、开车安全吗。两档都正确判断了明天南京无明显雾情，但对雾天驾驶
安全的论证方式截然不同。

\begin{reportsnippet}{ablated}{decision\_\bk{}006\_\bk{}fog · LLM-as-judge 专业性 3.0 / 5.0}
明天南京没有雾。凌晨湿度最高约 76\%，属于偏高但远未达到起雾条件，通常
需大于等于 90\% 以上且温差大才易形成辐射雾。

总体判断：明天开车是安全的。无雾、无雨，能见度正常；无天气预警生效；
风力不大，上午东南风 1 至 3 级，下午南风 3 至 4 级。
\end{reportsnippet}

\begin{reportsnippet}{full}{decision\_\bk{}006\_\bk{}fog · LLM-as-judge 专业性 5.0 / 5.0}
明天南京无明显雾情。全天无降水，凌晨湿度虽偏高（最高 76\%），但天气
状况为晴，不具备成雾条件；当前无大雾预警信息。

依据 GB/T 27964-2011《雾的预报等级》（中国气象局）：
轻雾对应能见度 1 至 10 km，影响较小，开示廓灯或近光灯并降速即可；
雾对应能见度 500 至 1000 m，需开近光灯加前后雾灯，保持车距；
浓雾对应能见度 50 至 500 m，触发大雾黄色预警，道路与机场显著受影响；
强浓雾对应能见度小于 50 m，触发红色预警，严重威胁交通安全。

明天南京预报为晴到阴，能见度远高于轻雾阈值（1 km），不存在雾天驾驶
风险。
\end{reportsnippet}

\texttt{ablated} 档给出的"通常需大于等于 90\% 才易形成辐射雾"是合理的
气象常识，但既无标准号支撑，也无能见度对照表帮助用户判断不同雾况下的
具体驾驶动作。\texttt{full} 档直接给出 GB/T 27964-2011《雾的预报等级》
4 档全表，把抽象的雾况判断转化为可对照的能见度区间与具体驾驶建议，
构成了完整可溯源的专业表达。LLM-as-judge 在专业性维度上对该条 ablated
档给 3.0 分而对 full 档给 5.0 分，2.0 分的差距正是双通路 RAG 在权威
分级表上的不可替代价值的直接量化。

第二个对照用例 \texttt{decision\_\bk{}001\_\bk{}outdoor} 是用户查询明天武汉
适合户外活动吗。两档都识别到明天夜间有大雨这一关键事实，但对大雨这一
分级的支撑力度差异显著。

\begin{reportsnippet}{ablated}{decision\_\bk{}001\_\bk{}outdoor · 关键事实 5/5、专业性 4/5}
明天武汉白天适合户外活动，建议安排在上午或下午 4 点以后，避开正午高温；
傍晚起需关注降雨，夜间大雨不宜外出。

注意事项：高温防暑（最高温达 32 °C，湿度高达 91\%，体感闷热）；夜间
大雨，不建议安排傍晚及夜间的户外活动，如需外出请携带雨具。
\end{reportsnippet}

\begin{reportsnippet}{full}{decision\_\bk{}001\_\bk{}outdoor · 关键事实 5/5、专业性 5/5}
明天武汉白天适合户外活动，但傍晚到夜间不适合。

夜间不适宜：傍晚后有大雨，根据《降水量等级》（GB/T 28592-2012），
大雨指 24 小时降水量 25.0 至 49.9 毫米，雨势猛烈，能见度明显下降，
地面易积水，不适合户外活动。

最佳时段为上午 6 至 12 点（晴到多云、温度适宜），次佳时段为下午
12 至 16 点（阴天但无雨，注意防晒），避免时段为傍晚 18 点后（有大雨，
建议室内活动）。
\end{reportsnippet}

\texttt{ablated} 档把大雨当作不证自明的常识词使用，回答语气自信但缺少
分级阈值依据。\texttt{full} 档则直接调用通路 B 给出 GB/T 28592-2012
中大雨档的精确区间 25.0 至 49.9 毫米与雨势猛烈、能见度明显下降、地面
易积水三类后果描述，把日常用语大雨转化为可量化、可决策的工程化表达。
LLM-as-judge 在专业性维度对 full 档给出满分 5.0；在 30 条端到端用例的
\texttt{decision\_\bk{}support} 类 6 条对照中，\texttt{ablated} 档通过率
仅 50\%，\texttt{full} 档为 100\%，差距 50 个百分点的端到端表现差距
即源自这一类报告生成阶段的引用纪律差异。

\subsubsection*{4.3.4 \quad 报告质量的端到端验证概要}

为量化上述报告生成机制的实际价值，本文采用 LLM-as-judge 4 维度评分
作为报告质量的代理指标，每条用例由独立大语言模型按事实性、完整性、
专业性、流畅性四个 0 至 5 区间维度打分。在 30 条覆盖六类场景的端到端
评测集上，两档对照结果的核心数据见 \reftab{tab:report-judge-summary}。

\begin{table}[!ht]
    \centering
    \caption{两档报告质量 LLM-as-judge 评分摘要（30 条端到端用例）}
    \label{tab:report-judge-summary}
    \song\wuhao
    \begin{tabular}{lccc}
        \toprule
        维度 & ablated & full & 差距 \\
        \midrule
        事实性 & 4.07 & 3.90 & $-0.17$ \\
        完整性 & 4.93 & 4.87 & $-0.06$ \\
        专业性 & 3.10 & 3.47 & $+0.37$ \\
        流畅性 & 4.97 & 5.00 & $+0.03$ \\
        综合 & 4.27 & 4.31 & $+0.04$ \\
        \bottomrule
    \end{tabular}
\end{table}

四个维度的差距分布揭示了报告生成阶段一个值得关注的现象。事实性、
完整性、流畅性三档接近持平甚至 \texttt{ablated} 档略胜，差距全部集中
在专业性维度的 $+0.37$ 个分差。这一结果可以从两个角度解读。第一，
大语言模型本身在表达流畅性与结构完整性上有较强的语料先验，无 RAG 时
也能写出看起来合理的报告。第二，当评估聚焦在分级、引用、阈值依据这
一专业维度时，双通路 RAG 的价值才真正显现出来。仅看综合分会得到 RAG
价值有限的错误印象，而拆分到专业性这一更细的维度才能反映真实差距。
完整的端到端评测设计、断言定义与按场景类别拆分的详细数据见
\ref{sec:e2e-eval} 节。

值得注意的是 \texttt{full} 档自身在专业性维度上的得分 3.47 仍未达到
满分。逐条分析失败用例可发现，4 条 \texttt{full} 档未通过的用例共同特征
是大语言模型在 RAG 结果中确实拿到了相关分级阈值，但在最终报告生成时
省略了 GB/T 编号。这一现象在 \ref{sec:semantic-bridge} 节描述为通路 B
单测层面的引用真实性已达到 100\% 但端到端集成后衰减到 26.7\%。这种从
模块层到集成层的指标衰减是单模块评测无法捕捉的盲区，构成
\ref{sec:future-work} 节关于
显式引用强制注入与生成后校验的具体改进方向。

% --- 4.3 节末段过渡（4 章整体小结迁出至独立 \subsection 4.4）---
报告生成阶段的设计与实证至此告一段落。从 \ref{sec:semantic-bridge} 节
双通路 RAG 算法到本节自适应报告模板，三组核心组件在数据流向上依次完成
数值到语义、数值到统计结果、多源信息到自然语言报告的串接，已具备进入
本章小结收口的条件，详见 \ref{sec:chap4-summary} 节。

